% ============================================================
% sections/02_problema.tex — Definição do Problema
% ============================================================

\section{Definição do Problema}

O ensino tradicional da tabela periódica baseia-se predominantemente em métodos
expositivos e mnemônicos, que exigem do estudante a memorização de um volume
considerável de informações desconexas do seu cotidiano. Essa abordagem frequentemente
resulta em aprendizado superficial e de curta duração \cite{gee2003videogames}.

O problema central que este trabalho busca endereçar pode ser formulado da
seguinte forma: \textit{como tornar o aprendizado dos elementos químicos da
tabela periódica mais engajador, significativo e eficaz por meio de um jogo
digital?}

% TODO: Adicionar referências a estudos sobre dificuldades no ensino de química
% TODO: Incluir dados de reprovação ou baixo desempenho em química no ensino médio

\subsection{Contexto Educacional}

Pesquisas recentes apontam que metodologias ativas de ensino, que colocam o
estudante como protagonista do processo de aprendizagem, produzem resultados
significativamente superiores em termos de retenção e compreensão do conteúdo
\cite{mayer2019serious}.

% Exemplo de equação (cálculo de eficiência de aprendizado — hipotético)
A eficiência de aprendizado $\eta$ pode ser modelada de forma simplificada como:

\begin{equation}
    \eta = \frac{C_r}{C_t} \times 100\%
    \label{eq:eficiencia}
\end{equation}

\noindent onde $C_r$ representa o conteúdo retido após determinado período e
$C_t$ o conteúdo total ensinado. A Eq.~\eqref{eq:eficiencia} ilustra que
metodologias que aumentam $C_r$ sem reduzir $C_t$ são intrinsecamente mais
eficientes.

\subsection{Lacuna Identificada}

Embora existam soluções digitais para o ensino de química, a maioria consiste
em simuladores ou quizzes que não exploram plenamente as mecânicas de engajamento
consolidadas pela indústria de jogos. O \textit{Alchemon} propõe preencher essa
lacuna ao incorporar sistemas de progressão, coleção e batalha diretamente
atrelados às propriedades dos elementos químicos.

% TODO: Realizar revisão sistemática de aplicativos existentes na área
% TODO: Comparar funcionalidades com solução proposta em tabela
